% ============================================================================
\section{Reproduction}
\label{app:repro}
% ============================================================================

From a clean checkout:

\begin{verbatim}
python -m venv .venv && . .venv/bin/activate     # Windows: .venv\Scripts\activate
python -m pip install -r requirements.txt
python -m pytest -q tests/                       # 370 tests
bash scripts/run_all.sh                          # Windows: scripts\run_all.ps1
\end{verbatim}

Individual campaigns, each of which writes \code{results/<id>/run\_record.json} with the exact
command, seeds, environment, duration and exit status:

\begin{verbatim}
python experiments/e1_welch_floor.py
python experiments/e2_threshold_recovery.py
python experiments/e3_linear_energy.py
python experiments/e4_optimized_codes.py
python experiments/e5_trained_toy.py
python experiments/e6_threshold_scaling.py
\end{verbatim}

Adding \code{--smoke} to any of them runs a reduced configuration end to end in seconds.
Figures, tables and the generated numbers are rebuilt with
\code{scripts/make\_figures.py}, \code{scripts/make\_tables.py} and
\code{scripts/make\_numbers.py}; \code{scripts/check\_manuscript\_numbers.py} verifies that
every measured number in this manuscript matches the saved results and exits non-zero
otherwise. E6 writes results per width and per chunk and resumes an interrupted campaign
without recomputing completed trials.

% ============================================================================
\section{Verification of imported statements}
\label{app:imported}
% ============================================================================

Every statement imported from the two source frameworks was checked against the published
version. Table~\ref{tab:imported} records the audit.

\begin{table}[h]
\centering
\caption{Audit of imported statements.}
\label{tab:imported}
\small
\begin{tabular}{p{3.1cm}p{4.6cm}p{4.6cm}}
\toprule
source & as printed there & as used here \\
\midrule
\cite[Thm.~11]{hanni2024superposition} &
targeted superpositional AND, precision $\varepsilon_{\mathrm{in}}\to\varepsilon_{\mathrm{out}}$
with $\varepsilon_{\mathrm{out}}=\Otilde(\sqrt{s^2/d})$ &
Lemma~\ref{lem:comp}, $\ecomp(d)=\Otilde(d^{-1/2})$ at $s=O(1)$ \\
\addlinespace[2pt]
\cite[Thm.~14]{hanni2024superposition} &
$\varepsilon^{(1)}=\Otilde(\max\{s\mu,\sqrt{s\varepsilon},\sqrt{s/d}\})$,
$\mu^{(1)}=\Otilde(\max\{\sqrt{1/d},\mu\})$ &
quoted in Sec.~\ref{sec:application} only to note it does not by itself preserve the
$\Otilde(d^{-1/2})$ invariant \\
\addlinespace[2pt]
\cite[Cor.~15]{hanni2024superposition} &
relaxes the near-sphere hypothesis of Thm.~14 at the cost of depth $1\to3$ &
not used for a capacity claim \\
\addlinespace[2pt]
\cite[Thm.~21]{hanni2024superposition} &
tolerance $\varepsilon<K\,d^{1/4}/(m^{1/2}s^{1/4})$, output
$\varepsilon^{(1)}=O(\log(d)\sqrt{s}/\sqrt d)$, $\mu^{(1)}=\Otilde(\sqrt s/\sqrt d)$ &
Lemma~\ref{lem:corr} with $m$ in the role of $F$ \\
\addlinespace[2pt]
\cite[Thm.~8]{hanni2024superposition} &
width $d=\Otilde(m^{2/3}s^2)$, depth $2L$ &
Proposition~\ref{prop:hanni}, inverted at $s=O(1)$ \\
\addlinespace[2pt]
\cite{adler2024superposition} &
$\Omega(\sqrt{m'}\log m')$ neurons and $\Omega(m'\log m')$ parameters; construction with
$O(\sqrt{m'}\log m')$ neurons; at most $O(n^2/\log n)$ features with $n$ neurons &
Sec.~\ref{sec:application}; the parameter-to-neuron conversion is flagged as an architectural
assumption of that model \\
\addlinespace[2pt]
\cite[Thm.~2.3]{strohmer2003grassmannian} &
$\max_{i\neq j}|\langle\phi_i,\phi_j\rangle|\ge\sqrt{(F-d)/(d(F-1))}$ for unit-norm systems &
the Gram case of Theorem~\ref{thm:floor}; the reduction is noted in Sec.~\ref{sec:frames} \\
\bottomrule
\end{tabular}
\end{table}

% ============================================================================
\section{Auxiliary proofs}
\label{app:proofs}
% ============================================================================

\subsection{Sub-Gaussian interference}

Let $u,u_1,\dots,u_m$ be independent uniform unit vectors in $\R^d$. Conditionally on $u$, each
$X_j=\langle u,u_j\rangle$ is distributed as the first coordinate of a uniform point on
$S^{d-1}$, hence has mean zero and is sub-Gaussian with variance proxy $C_0/d$ for a universal
$C_0$~\cite[Thm.~3.4.6]{vershynin2018hdp}. The $X_j$ are conditionally independent, so their
sum is sub-Gaussian with proxy $C_0m/d$ and
$\Pr\{|\sum_j X_j|>t\}\le2\exp(-c_0dt^2/m)$ with $c_0=1/(2C_0)$. This is the estimate used in
Theorem~\ref{thm:randomsupport}.

\subsection{The Bernoulli variant of the energy floor}

Under $b_i\sim\mathrm{Bernoulli}(p)$ independently,
$\E[bb^{\top}]=p(1-p)I_F+p^2\mathbf 1\mathbf 1^{\top}$, so
$\E_b\norm{Ab}^2=p(1-p)\norm{A}_F^2+p^2\norm{A\mathbf 1}^2\ge p(1-p)F(F-d)/d$ by
Theorem~\ref{thm:floor}. With $p=s/F$ and $s\le F/2$ this gives
$\frac1F\E_b\norm{Ab}^2\ge\frac{s(F-d)}{2dF}=\Omega(s/d)$, matching
Proposition~\ref{prop:uniform-energy} up to the $(F-1)$ versus $F$ denominator. Both variants
are computed by the released code and both are reported in E3.

% ============================================================================
\section{The nuclear-norm inequality}
\label{app:nuclear}
% ============================================================================

For $M\in\R^{F\times F}$ with singular value decomposition
$M=\sum_{k=1}^{r}\sigma_k u_k v_k^{\top}$,
\[
  |\operatorname{tr}(M)|=\Big|\sum_k\sigma_k\,u_k^{\top}v_k\Big|
  \le\sum_k\sigma_k|u_k^{\top}v_k|\le\sum_k\sigma_k=\nucnorm{M},
\]
and by Cauchy--Schwarz on $(\sigma_1,\dots,\sigma_r)$,
$\nucnorm{M}^2\le r\sum_k\sigma_k^2=r\norm{M}_F^2$. With $r\le d$ this gives
$|\operatorname{tr}(M)|^2\le d\norm{M}_F^2$, as used in Theorem~\ref{thm:floor}. Note that the
identity $\operatorname{tr}(M)=\sum_k\sigma_k$ holds only for positive semidefinite matrices and
is \emph{not} used.

% ============================================================================
\section{Reference scales}
\label{app:scales}
% ============================================================================

The scales appearing in this literature count different objects and should not be read as a
capacity ordering. Writing $F_{\mathrm{CS}}(d,\alpha)=d\,g(\alpha)$ with
$g(\alpha)=1/((1-\alpha)\ln(1/(1-\alpha)))$ for compressed-sensing-style linearly decodable
storage, $F^{\mathrm{H}}_{\mathrm{cert}}(d)=\Thetatilde(d^{3/2})$ for the H\"{a}nni template
certificate, $L_{\mathrm{AS}}(d)=d^2/\log^2 d$ and $U_{\mathrm{AS}}(d)=d^2/\log d$ for the two
sides of the Adler--Shavit bracket, and $N_{\mathrm{JL}}(d;\varepsilon)=\exp(\Theta(d
\varepsilon^2))$ for Johnson--Lindenstrauss packing~\cite{JL1984}, one has
\[
  F_{\mathrm{CS}}\ll F^{\mathrm{H}}_{\mathrm{cert}}\ll L_{\mathrm{AS}}
  \lesssim\FrecAS\lesssim U_{\mathrm{AS}}\ll N_{\mathrm{JL}}
\]
as $d\to\infty$. This does not mean that recursive computation beats storage:
$N_{\mathrm{JL}}$ counts passively packable states, $F_{\mathrm{CS}}$ counts linearly decodable
features through a compressed bottleneck, and the middle two count recursively computable
Boolean features under two different published frameworks. Michaud et
al.~\cite{michaud2025manifolds} give a complementary reason why an observed dictionary size may
fall below the template scale: if feature manifolds consume capacity, the relevant comparison
is the effective dimension rather than the raw feature count.

% ============================================================================
\section{Sparse-autoencoder context}
\label{app:sae}
% ============================================================================

As background only, \citet{borobia2026pruning} train TopK sparse
autoencoders ($k=64$, expansion $8d$) on three language models and observes no dead features at
$F=8d$, placing those dictionaries well below both the $d^{3/2}$ template scale and the
near-quadratic Adler--Shavit scale. Sparse-autoencoder dictionary size is a reconstructive
quantity and is not a measurement of recursive computation capacity; these numbers are not used
in any proof or experiment of the present paper, and are mentioned only to indicate where
current empirical dictionaries sit relative to the scales of Appendix~\ref{app:scales}.
Liu et al.~\cite{liu2025scaling} and Elhage et al.~\cite{elhage2022superposition} provide
further context on why strong superposition regimes are of interest.
